% GENERATED FILE. Edit canonical lessons, then run npm run generate:books.
% canonical-source-hash: fnv1a64:b052b9171a15399e
% canonical-lessons: BN-C170-akar, BN-C170-tukro, BN-C170-lain, BN-C170-michhil, BN-C170-kol

\chapter{Shape, Piece, Queue, Procession, Tap}
\label{ch:bn-shape-piece-queue-procession-tap}
\hlchaptermodality{170}

\begin{chapteropening}
\textbf{By the end of this chapter:} \emph{I can say a shape, a piece, a queue, a procession and a tap.}

Complete the last lesson of chapter 170: I can say a shape, a piece, a queue, a procession and a tap.
\end{chapteropening}

\section[\texorpdfstring{\bn{আকার} (ākār)}{আকার (ākār)}]{\bn{আকার} (ākār) — a shape, a size}

\label{lesson:BN-C170-akar}



\begin{quote}
Before the new one: say the Bengali for anger, then the Bengali for peace.
\end{quote}

\subsection*{You'll want to know: \bn{আকার}}
\textbf{\bn{আকার}} — \emph{ākār} — \textquotedblleft{}a shape, a size\textquotedblright{}.

\begin{cousinweb}
\textbf{\bn{আকার}} (\emph{ākār}) is Sanskrit \emph{ākāra}, \textquotedblleft{}shape\textquotedblright{}, taken over whole as a learned loan, a tatsama (CDIAL 1004).
\end{cousinweb}

\begin{grammarlens}[title={What you've built}]
One more word for this chapter.
\end{grammarlens}

\subsection*{Your turn}
\begin{itemize}
\raggedright
  \item \emph{Say it:} \emph{ākār}
  \item \emph{Say it:} \emph{ākār}, once more
  \item \emph{Say it:} say \emph{ākār}
  \item \emph{Recall:} say the Bengali for anger, then the Bengali for peace, then say \emph{ākār} again
\end{itemize}

\begin{tcolorbox}[breakable,colback=teal!4,colframe=teal!35!black,title={Before you move on}]
What does \bn{আকার} mean? (\textquotedblleft{}A shape, a size\textquotedblright{}.) Say it once more.
\end{tcolorbox}
\section[\texorpdfstring{\bn{টুকরো} (ṭukro)}{টুকরো (ṭukro)}]{\bn{টুকরো} (ṭukro) — a piece}

\label{lesson:BN-C170-tukro}



\begin{quote}
Before the new one: say the Bengali for peace, then the Bengali for a shape.
\end{quote}

\subsection*{You'll want to know: \bn{টুকরো}}
\textbf{\bn{টুকরো}} — \emph{ṭukro} — \textquotedblleft{}a piece\textquotedblright{}.

\begin{cousinweb}
\textbf{\bn{টুকরো}} (\emph{ṭukro}) is a tadbhava, a word inherited and reshaped by sound change. It comes from \emph{*ṭukka}, \textquotedblleft{}piece\textquotedblright{}, a form CDIAL reconstructs from its descendants (CDIAL 5466). CDIAL lists Bengali \emph{ṭukrā}, \textquotedblleft{}small piece\textquotedblright{}, made with an extra \emph{-r-}.
\end{cousinweb}

\begin{grammarlens}[title={What you've built}]
One more word for this chapter.
\end{grammarlens}

\subsection*{Your turn}
\begin{itemize}
\raggedright
  \item \emph{Say it:} \emph{ṭukro}
  \item \emph{Say it:} \emph{ṭukro}, once more
  \item \emph{Say it:} say \emph{ṭukro}
  \item \emph{Recall:} say the Bengali for peace, then the Bengali for a shape, then say \emph{ṭukro} again
\end{itemize}

\begin{tcolorbox}[breakable,colback=teal!4,colframe=teal!35!black,title={Before you move on}]
What does \bn{টুকরো} mean? (\textquotedblleft{}A piece\textquotedblright{}.) Say it once more.
\end{tcolorbox}
\section[\texorpdfstring{\bn{লাইন} (lāin)}{লাইন (lāin)}]{\bn{লাইন} (lāin) — a queue, a line}

\label{lesson:BN-C170-lain}



\begin{quote}
Before the new one: say the Bengali for a shape, then the Bengali for a piece.
\end{quote}

\subsection*{You'll want to know: \bn{লাইন}}
\textbf{\bn{লাইন}} — \emph{lāin} — \textquotedblleft{}a queue, a line\textquotedblright{}.

\begin{cousinweb}
\textbf{\bn{লাইন}} (\emph{lāin}) is a loan from English \emph{line}.
\end{cousinweb}

\begin{grammarlens}[title={What you've built}]
One more word for this chapter.
\end{grammarlens}

\subsection*{Your turn}
\begin{itemize}
\raggedright
  \item \emph{Say it:} \emph{lāin}
  \item \emph{Say it:} \emph{lāin}, once more
  \item \emph{Say it:} say \emph{lāin}
  \item \emph{Recall:} say the Bengali for a shape, then the Bengali for a piece, then say \emph{lāin} again
\end{itemize}

\begin{tcolorbox}[breakable,colback=teal!4,colframe=teal!35!black,title={Before you move on}]
What does \bn{লাইন} mean? (\textquotedblleft{}A queue, a line\textquotedblright{}.) Say it once more.
\end{tcolorbox}
\section[\texorpdfstring{\bn{মিছিল} (michhil)}{মিছিল (michhil)}]{\bn{মিছিল} (michhil) — a procession}

\label{lesson:BN-C170-michhil}



\begin{quote}
Before the new one: say the Bengali for a piece, then the Bengali for a queue.
\end{quote}

\subsection*{You'll want to know: \bn{মিছিল}}
\textbf{\bn{মিছিল}} — \emph{michhil} — \textquotedblleft{}a procession\textquotedblright{}.

\begin{cousinweb}
The origin of \textbf{\bn{মিছিল}} (\emph{michhil}) is not settled: CDIAL has no entry for it, and a search by form and by the gloss \textquotedblleft{}procession\textquotedblright{} finds nothing close.
\end{cousinweb}

\begin{grammarlens}[title={What you've built}]
One more word for this chapter.
\end{grammarlens}

\subsection*{Your turn}
\begin{itemize}
\raggedright
  \item \emph{Say it:} \emph{michhil}
  \item \emph{Say it:} \emph{michhil}, once more
  \item \emph{Say it:} say \emph{michhil}
  \item \emph{Recall:} say the Bengali for a piece, then the Bengali for a queue, then say \emph{michhil} again
\end{itemize}

\begin{tcolorbox}[breakable,colback=teal!4,colframe=teal!35!black,title={Before you move on}]
What does \bn{মিছিল} mean? (\textquotedblleft{}A procession\textquotedblright{}.) Say it once more.
\end{tcolorbox}
\section[\texorpdfstring{\bn{কল} (kôl)}{কল (kôl)}]{\bn{কল} (kôl) — a tap}

\label{lesson:BN-C170-kol}



\begin{quote}
Before the new one: say the Bengali for a queue, then the Bengali for a procession.
\end{quote}

\subsection*{You'll want to know: \bn{কল}}
\textbf{\bn{কল}} — \emph{kôl} — \textquotedblleft{}a tap\textquotedblright{}.

\begin{cousinweb}
CDIAL lists Bengali \emph{kal}, \textquotedblleft{}lock of a gun, machine\textquotedblright{}, under Sanskrit \emph{kalā}, \textquotedblleft{}any practical art\textquotedblright{}, a word it takes from Dravidian (CDIAL 2927). A tap is a \emph{kôl} in the sense \textquotedblleft{}a device\textquotedblright{}.
\end{cousinweb}

\begin{grammarlens}[title={What you've built}]
One more word for this chapter.
\end{grammarlens}

\subsection*{Your turn}
\begin{itemize}
\raggedright
  \item \emph{Say it:} \emph{kôl}
  \item \emph{Say it:} \emph{kôl}, once more
  \item \emph{Say it:} say \emph{kôl}
  \item \emph{Recall:} say the Bengali for a queue, then the Bengali for a procession, then say \emph{kôl} again
\end{itemize}

\begin{tcolorbox}[breakable,colback=teal!4,colframe=teal!35!black,title={Before you move on}]
What does \bn{কল} mean? (\textquotedblleft{}A tap\textquotedblright{}.) Say it once more.
\end{tcolorbox}
